\documentclass[10pt,a4paper]{amsart}
\usepackage[margin=19mm]{geometry}
\usepackage{amsmath,amssymb,mathtools}
\usepackage[scr=rsfs,cal=euler]{mathalfa}
\usepackage{xfrac}
\usepackage[unicode,hidelinks,bookmarks=false]{hyperref}
\newcommand{\ie}{i.e.\ }
\newcommand{\cf}{cf.\ }
\newcommand{\resp}{respectively\ }
\newcommand{\Subsection}{\subsection}
\providecommand{\nobreakdash}{\nobreak\hbox{-}}
\setlength{\emergencystretch}{2em}
\multlinegap0pt
\usepackage{amssymb}
\usepackage{mathtools}
\usepackage{xcolor}
\newtheorem{lemma}{Lemma}
\newcommand{\jim}[1]{{ \oo{[Jim comment: #1}]}}
\newcommand{\matip}[2]{#1 \bullet #2}
\newcommand{\oo}[1]{{\textcolor{blue}{#1}}}
\newcommand{\ooc}[1]{{ \oo{[OG comment: #1}]}}
\title{Sparse Cuts for the Positive Semidefinite Cone}
\author{Oktay G\"unl\"uk$^a$ \and Paul J\"unger$^b$ \and {Jeff Linderoth}$^c$ \and Andrea Lodi$^b$ \and James Luedtke$^c$}
\date{Source: arXiv:2603.09864v2 (CC BY 4.0)}
\begin{document}
\maketitle

\noindent\emph{Source and licence.} Sparse Cuts for the Positive Semidefinite Cone. Version arXiv:2603.09864v2 (CC BY 4.0), distributed by arXiv under the Creative Commons Attribution 4.0 International licence, http://creativecommons.org/licenses/by/4.0/, as recorded in the arXiv licence metadata for this exact version and shown on the abstract page https://arxiv.org/abs/2603.09864v2. Full licence notice: \texttt{licenses/arxiv-2603.09864-CC-BY-4.0.txt}.\par\smallskip

\noindent\textit{Editorial scope.} Counted unit: the article's lemma lem:dnnform (source0.tex lines 548--557). The article's complete original proof of it is delivered with it (source0.tex lines 558--582, one \texttt{proof} environment of 25 lines). No mathematics is written, repaired or re-derived: the statement and the proof are the article's own lines, in the article's own notation, and no retained line is substituted or re-flowed.  Retained uncounted as the source-local context the proof needs: lines 530--533; lines 543--546.  Nothing was omitted from any retained span, so the package carries no omission record.  No retained line contains a citation, so the wrapper carries no bibliography and no bibliography entry.  No retained line contains a figure environment, an image-inclusion command or drawing code, so no image is delivered and the package declares no asset dependency.  Editorial formatting only, in addition: an AMS \texttt{amsart} preamble that restates the article's own declarations and macros used by the retained lines, namely the environments \texttt{lemma} on separate counters, together with the standard packages those lines need.  The retained environments print in this document as Lemma 1 in order of appearance, whereas the article prints the same items with the numbering of its own full document.  The complete native source of the article as distributed in the arXiv e-print for this exact version is bundled unchanged as \texttt{sources/196093/source0.tex}.
	\begin{equation}
		\label{eq:sdpsep}
		\min_{{C\in \mathcal S}} \Bigl\{ \matip{C}{\widehat{Y}_{\mathrm{LP}}}  : C \succeq 0, \ C_{ij} = 0 \ \forall(i,j) \notin E, \ \| \mathrm{diag}(C)\|_1 \leq 1 \Bigr\},
	\end{equation}

	\begin{equation}
		\label{eq:dnnsep-init}
		\min_{C \in \mathcal S} \Bigl\{ \matip{C}{\widehat{Y}_{\mathrm{LP}}}  : C \in \mathcal{D}^*, \ C_{ij} = 0 \ \forall(i,j) \notin E, \ \| \mathrm{diag}(C) \|_1  \leq 1 \Bigr\}.
	\end{equation}

	\begin{lemma}
		\label{lem:dnnform}
		Problem \eqref{eq:dnnsep-init} is equivalent to problem 
		
		%\ooc{any objections to replacing $C \succeq 0$ below to $C\in \mathcal{D}^*$ ?}\jim{This change would defeat the purpose, which is to demonstrate the problem is an SDP.}
		\begin{equation}
			\label{eq:dnnsep}
			\min_{{C\in \mathcal S}} \Bigl\{ \matip{C}{\widehat{Y}_{\mathrm{LP}}} : C \succeq 0, \ C_{ij} \leq 0 \ \forall(i,j) \notin E, \ \| \mathrm{diag}(C)\|_1 \leq 1 \Bigr\} .
		\end{equation}
	\end{lemma}

	\begin{proof}
		By definition, the following semidefinite optimization problem is a restatement of %the problem 
		\eqref{eq:dnnsep-init}, where we use $A$ in place of $C$ as the matrix defining the cut:
		\begin{align*}
			\min_{{A\in \mathcal S}} \ & \matip{A}{\widehat{Y}_{\mathrm{LP}}}  \\
			\text{s.t. } & A = C +D \\
			& D \geq 0 \\
			& C \succeq 0\\
			& A_{ij} = 0 \quad \text{for all } (i,j) \notin E\\
			& \| \mathrm{diag}(A) \|_1 \leq 1 .
		\end{align*}
		Since $\widehat{Y}_{\mathrm{LP}} \geq 0$,  there always exists an optimal solution with $D_{ij} = 0$ and $C_{ij} = A_{ij}$  for all $(i,j) \in E$.
		This also implies that $\mathrm{diag}(A)=\mathrm{diag}(C)$. Also note that $A_{ij} = 0$ implies that $C_{ij} = -D_{ij}\le 0$ for all $(i,j) \notin E$. 
		Recalling that $(j,j) \in E$ for $j=0,\ldots,n$, we can then eliminate the $A$ and $D$ variables to  simplify the separation problem to yield \eqref{eq:dnnsep}.
		%\begin{subequations}
		%\label{eq:projsdp-dnn}
		%\begin{align}
		%\min \ & \langle C_E,\hat{Z} \rangle \\
		%\text{s.t. } & C \succeq 0 \\
		%   & C_{ij} \leq 0 \quad \text{for all } (i,j) \notin E \label{eq:projsdp-dnn-sparsity}\\
		%   & \| \mathrm{diag}(C) \|_1 \leq 1.
		%\end{align}
		%\end{subequations} 
		%Comparing this formulation to \eqref{eq:sdpsep} yields the claim. 
	\end{proof}

\end{document}
